\documentclass[12pt]{exam}

\usepackage[T1]{fontenc}
\usepackage[utf8]{inputenc}
\usepackage[lithuanian]{babel}
\usepackage{lmodern}
\usepackage{graphicx}
\usepackage{xcolor}
\usepackage{tikz}
\usepackage[
    a4paper,
    left=1cm,
    right=1cm,
    top=2.5cm,
    bottom=1cm
]{geometry}
\usepackage{amsmath,amssymb}
\usepackage{tasks}

% Švelnus pilkas rėmelis mokinio įrašui
\newcommand{\irasolaukas}[1]{%
    \begingroup
    \setlength{\fboxrule}{0.3pt}%
    \setlength{\fboxsep}{3pt}%
    \fcolorbox{black!35}{white}{%
        \rule{0pt}{0.45cm}\hspace{#1}%
    }%
    \endgroup
}

% Sprendimo laukas:
% argumentas nurodo linijų skaičių;
% tarpas tarp linijų – 0.7 cm.
\newcommand{\sprendimas}[1]{%
    \par\smallskip
    \noindent\textit{Sprendimas:}
    \par\nobreak
    \vspace{0.1cm}
    \noindent
    \begin{tikzpicture}
        \path[use as bounding box]
            (0,0) rectangle (\linewidth,{-0.7*#1});
        \foreach \i in {1,...,#1} {
            \draw[
                draw=black!25,
                line width=0.3pt
            ]
                (0,{-0.7*\i})
                -- (\linewidth,{-0.7*\i});
        }
    \end{tikzpicture}%
    \par\smallskip
}

% Atsakymo laukelis
\newcommand{\atsakymas}[1][9]{%
    \par\smallskip
    \noindent\textit{Atsakymas:}
    \framebox[#1cm]{\rule{0pt}{0.5cm}}
    \par
}

% Antraštė
\pagestyle{headandfoot}
\runningheadrule

\firstpageheader
{%
    \small
    \begin{tabular}{@{}l@{\hspace{0.2cm}}l@{}}
        Mokinio kodas: & \irasolaukas{2.3cm} \\[0.15cm]
        Klasė:        & \irasolaukas{2.3cm}
    \end{tabular}%
}
{Kontrolinis darbas:\\Kvadratinės lygtys}
{Prometėjo variantas}

\runningheader{Kvadratinės lygtys}{}{}
\firstpagefooter{}{}{}
\runningfooter{}{}{}

% Taškų žymėjimas
\pointsdroppedatright
\bracketedpoints
\setlength{\rightpointsmargin}{1.6cm}
\renewcommand{\points}{taškai}
\renewcommand{\pointname}{Taškai:}
\marginpointname{\ taškai}

\begin{document}

\begin{questions}

\question[3]
Išspręskite lygtį $(2x-1)^2=(x+2)^2$.\droppoints
\sprendimas{2}
\atsakymas[3]

\question[3]
Lygties $x^2-2x-1=0$ sprendiniai yra $x_1$ ir $x_2$. Raskite lygties sprendinių kubų sumą. \droppoints
\sprendimas{2}
\atsakymas[5]

\question[6]
Išskaidykite dauginamaisiais:\droppoints
\begin{tasks}[label-width=1.2em](3)
\task $-2x^2+5x-3$
\task $2x^2-xy+6x-3y$
\task $x^3-8$
\end{tasks}
\sprendimas{3}
\atsakymas[9]

\question[4]
Du mokiniai turėjo išspręsti lygtį $x^2+bx+c=0$. Pirmasis neteisingai nurašė tik koeficientą $c$, tačiau savo užrašytą lygtį išsprendė teisingai ir gavo sprendinius
$2$ ir $5$. Antrasis neteisingai nurašė tik koeficientą $b$, tačiau savo užrašytą lygtį išsprendė teisingai ir gavo sprendinius $-1$ ir $6$. Raskite $b$ ir $c$. 

\droppoints
\sprendimas{3}
\atsakymas[4]

\question[5]
Išspręskite lygtį $2(2x^2-6x+1)(x^2-3x-2{,}5)=27$.\droppoints
\sprendimas{4}
\atsakymas[4]

\question[5]
Duotas kvadratinis trinaris  $P(x) = x^2-3x-1$. Apskaičiuokite $a^3-2a^2-4a+7$ reikšmę, kur $a$ yra viena iš šio trinario šaknų. \droppoints
\sprendimas{6}
\atsakymas[3]

\question[7]
Raskite visas realiąsias $m$ reikšmes, su kuriomis lygtis $x^2-(m+1)x+2m-2=0$ turi du skirtingus teigiamus sprendinius, kurių vienas tris kartus didesnis už kitą.\droppoints
\sprendimas{7}
\atsakymas[4]

\question[7]
Lygties $x^2-4x+1=0$ sprendiniai yra $x_1$ ir $x_2$. Neskaičiuodami jų, sudarykite redukuotą kvadratinę lygtį, kurios sprendiniai būtų $(x_1+1)^2$ ir $(x_2+1)^2$.\droppoints
\sprendimas{7}
\atsakymas[6]
\end{questions}

% Vertinimo lentelė puslapio apačioje
\vfill

\begin{center}
    \renewcommand{\arraystretch}{1.3}
    \begin{tabular}{|c|c|c|c|c|c|c|c|c|c|c|}
        \hline
        \textbf{Taškai}
        & 0--4 & 5--8 & 9--12 & 13--16 & 17--20
        & 21--24 & 25--29 & 30--33 & 34--37 & 38--40
        \\ \hline
        \textbf{Pažymys}
        & \textbf{1} & \textbf{2} & \textbf{3} & \textbf{4} & \textbf{5}
        & \textbf{6} & \textbf{7} & \textbf{8} & \textbf{9} & \textbf{10}
        \\ \hline
    \end{tabular}
\end{center}

\end{document}
